Reinforcement learning from execution feedback (RLEF) with binary rewards suffers from gradient starvation: at group size $G=4$, up to 69\% of GRPO training steps produce zero gradient because the model either passes or fails every completion in a group. We propose offline variance-guided selection --- profiling a frozen model once to compute per-problem binary execution reward variance $p_i(1-p_i)$, then selecting problems with the highest variance before training begins. On MBPP (374 problems) with DeepSeek-Coder-7B-Instruct, we find that 7.8\% of problems exhibit nonzero variance under constrained generation ($k=4$, \texttt{max\_new\_tokens=128}), and that top-50 variance-guided selection achieves $4.24\times$ higher mean variance than random-50 selection (Mann-Whitney $p=2.22\times10^{-6}$) --- stochastically dominating random selection across the full distribution, with profiling completing in $\sim$32 seconds via vLLM. However, GRPO training on the variance-selected subset produces no gradient signal across 50,000+ generation attempts (50--200 steps), because the model generates zero correct solutions under training conditions --- a cold-start precondition that renders data selection method irrelevant regardless of analytical validity. We identify generation parameter mismatch between profiling and training as the primary candidate root cause (among several alternatives) and provide a prescriptive resolution protocol --- proposed but not yet empirically validated. Our contribution is not a working method but a precisely characterized failure mode and its necessary precondition: the kind of knowledge that prevents others from repeating 50,000 wasted computation attempts. Our work establishes cold-start verification as a necessary precondition for RLEF data selection methods, and provides, to our knowledge, the first systematic per-problem empirical characterization of binary execution reward variance distribution on MBPP for a 7B code LLM.
